\chapter*{Conferencia XVII}
\addcontentsline{toc}{chapter}{LECTURA I - El timbre y el carácter del violín}
SEÑORES: Las razones filosóficas para la superficie interior permanente y perfectamente lisa del violín están claramente definidas y son fáciles de comprender. Tras largos años de observación sobre este tema, me parece innecesario exponer las razones para proteger las superficies interiores del violín de fuerzas desintegradoras como el calor y la humedad. También me parece innecesario exponer las razones para aumentar la potencia del sonido gracias a la superficie interior perfectamente lisa; pero, recordando mis propias dudas sobre estos puntos antes de mis observaciones, supongo que otros, que estén considerando estos puntos por primera vez, también puedan tener dudas similares.
dudas.
Al principio, mis propias dudas no me permitieron probar este experimento en violines valiosos; ni lo probé en violines valiosos hasta después de varios años de observación. En este momento, me pregunto por qué el mundo del violín debería temer más la protección de las superficies internas del violín que de las externas. Que tal temor existe en gran medida es bien sabido; pero las razones de tal temor no están del todo claras; algunos sostienen una objeción, otros otra. Quizás la objeción más universal se deba al temor de que la protección interna pueda ser desastrosa para la calidad del sonido. No se puede culpar al propietario de un violín con una calidad de sonido superior por albergar tal temor, y, dado que los violines de valor tonal superior son relativamente raros, si tales violines fueran tan comunes como los violines de calidad tonal inferior, entonces el caso sería diferente.
256

De nuevo, la situación sería diferente si los violinistas vivieran lo suficiente como para ver cómo se desgastan los violines nuevos bajo sus propios arcos. Así, al saber que un violín valioso está perdiendo calidad, el propietario adoptaría con gusto cualquier método razonable para prolongar su vida útil. Sin embargo, el desgaste de un violín en el transcurso de una sola vida es un hecho poco común.
Basándonos en cálculos de igual uso, no se necesita argumentar para determinar que, de dos violines construidos con las mismas secciones de troncos de madera dura y blanda, el que tenga menos madera está destinado a un período de utilidad más breve. Recuerdo vívidamente que cuando la rigidez de la caja de resonancia se reduce hasta el punto que produce la mayor amplificación del tono, a partir de entonces, una ligera pérdida de madera por desintegración disminuye la potencia tonal; y dicha pérdida, continua debido a la continua desintegración, puede arruinar un valor tonal superior en el transcurso de una sola vida. Por lo tanto, podemos esperar una disminución inmediata del valor tonal de los violines nuevos que poseen una potencia tonal superior en el momento de salir de las manos del constructor. Asimismo, y precisamente por la continua desintegración en las superficies interiores desprotegidas, podemos esperar una mayor potencia tonal de los violines nuevos que tienen un poco más de madera en el momento de salir de las manos del constructor. Es humano desear una superioridad inmediata.
en los valores tonales del violín.
La vida es demasiado corta para esperar una lenta desintegración.
Una vez poseyó valores tonales de violín superiores, 'tis 257

Pero es humano llorar su pérdida. Prevenir tal pérdida es sumamente sencillo.
Siento una punzada de pesar cada vez que pienso en esas joyas destrozadas esparcidas a lo largo del camino del violín; y, porque casi todos esos fragmentos, debido a la desintegración de sus superficies internas, podrían haberse conservado por un período indefinido.
Los detalles sobre la protección de la superficie interior del violín, según mi conocimiento, se han presentado en ocasiones anteriores; sin embargo, en dichas presentaciones, dudé en abordar los principios filosóficos que la sustentan. Mi vacilación se debía a que relativamente pocos violinistas se interesan por la filosofía relacionada con la producción o la conservación del sonido del violín. Es bastante seguro asumir que todos los violinistas buscan con avidez afirmaciones breves y categóricas de que tal o cual cosa, tal o cual método, produce infaliblemente ciertos resultados en el sonido del violín. He observado que, incluso hoy, después de que el violín haya sido objeto de atención durante 400 años, nadie puede formular reglas infalibles para la producción de los mejores valores tonales. Cuando llegue el día en que existan tales reglas infalibles, los violines con los mejores valores tonales serán tan comunes como lo son hoy los violines con valores tonales inferiores. El hecho es que fabricar violines con un sonido defectuoso es sumamente fácil y seguro; mientras que fabricar violines con un sonido impecable es sumamente difícil e incierto. Tal incertidumbre se debe a la imposibilidad de formular reglas que rijan todos los fenómenos en violación
tono lin.	En cuanto a la causa de tal impossi-258

En cuanto a la capacidad, existe un desacuerdo entre los filósofos de hace 50 años y los filósofos de hoy. Hoy, la mayoría sostiene la opinión de que tal imposibilidad se debe a la acción caprichosa de la madera de la caja de resonancia. Mi propia experiencia corrobora decididamente esta opinión. Creo que hoy, con madera de la mejor calidad para la caja de resonancia, se podrían producir violines con los mejores valores tonales indefinidamente. Es evidente que el constructor, incluso cuando dispone de la mejor madera, debe necesariamente dominar todos los demás factores que operan para modificar el tono del violín. En mi trabajo sobre el violín, ha sido un objetivo separar y clasificar todos los factores involucrados en la
producción y modificación del tono del violín de la siguiente manera:
Clase 1. Clase 2.	Factores constantes en los resultados tonales. Factores que no son constantes en los resultados tonales.
Los siguientes principios modificadores del tono, al ser constantes en sus resultados, se colocan en la Clase I, de la siguiente manera:	
1.	Alargar un agente productor de tono, otros
Si las dimensiones permanecen iguales, el tono disminuye.
2. Acortar un agente productor de tono, manteniendo las demás dimensiones iguales, eleva el tono.
3. Al aumentar el grosor de un agente productor de tono, manteniendo las demás dimensiones iguales, aumenta el tono.
4. Al disminuir el grosor de un agente productor de tono, manteniendo las demás dimensiones iguales, disminuye el tono.
5. El alargamiento de las columnas de aire perpendiculares y confinadas reduce el tono.
6. El acortamiento de las columnas de aire perpendiculares y confinadas eleva el tono.
259

7. La posición de las salidas del violín, alejada de los puntos de concentración de las ondas sonoras, disminuye el volumen del tono, la intensidad del tono y la calidad del sonido, que resulta ruidoso.
8. El aumento del área de salidas del violín eleva el tono, aumenta el volumen del sonido, disminuye la intensidad del sonido y acentúa la calidad del sonido ruidoso.
9. La disminución del área de salida del violín reduce el tono, disminuye el volumen del sonido, aumenta la intensidad del sonido y disminuye la calidad del sonido ruidoso.
10. Las superficies interiores rugosas y alfombradas del violín disminuyen el volumen del sonido, la intensidad del sonido y la calidad del sonido ruidoso.
11. Las superficies interiores perfectamente lisas del violín aumentan la potencia del tono y acentúan el ruido.
calidad del tono.
Los siguientes factores modificadores del tono, al no ser constantes en su acción, se colocan en la Clase II, de la siguiente manera:
1. El diapasón.
2. El puente.
3. El banco de resonancia.
4. El
correo.
5. La parte de atrás.
6. El.bar.
7. Los bloques.
8. Las costillas.
9. El
revestimientos.
Algunos de los modificadores de la Clase II se aproximan casi a una acción constante, o más bien a la falta de acción, como el diapasón, los bloques, las costillas y los revestimientos; pero, debido a que carecen de efectos constantes sobre el tono en cualquier grado, por lo tanto se colocan en esta clase. [En
el momento de hacer la lista en la Clase II, yo
260

He descubierto que el efecto que tienen las costillas, los forros y los bloques sobre el tono no ha recibido atención hasta ahora. Debido a esta omisión, creo que puede haber otras. Ofrezco las dificultades con las que trabajo como excusa. He observado que el beneficio que aportan las costillas, los forros y los bloques al tono consiste en evitar la vibración excesiva del violín. Dado que dicha vibración debilita el tono, el peso y la rigidez de estos elementos son la medida de su valor para el sonido; y, por lo tanto, esta medida refleja su falta de acción.
Se observa que todos los modificadores de tono de la clase I dependen completamente de la acción del aire; por lo tanto, sin lugar a dudas, su efecto constante sobre el sonido del violín se debe a que la acción del aire es una magnitud constante. Los primeros seis modificadores de tono de la clase I afectan únicamente a la cualidad tonal de la altura; el séptimo afecta a tres cualidades tonales; el octavo y el noveno, a cuatro cualidades tonales cada uno; el décimo y el undécimo, a tres cualidades tonales cada uno.
En este punto, es interesante notar la cantidad de cualidades tonales bajo el control absoluto del constructor. Encuentro que esta cantidad no es lo suficientemente grande como para ser halagadora. Enumerando las cualidades tonales que intervienen en el tono "rico" del violín, encuentro que su número es igual a 12, así: tono, volumen, intensidad, uniformidad, ausencia de armónicos disonantes o ruido, simpatía en concierto, respuesta a la presión del arco, tonos agradables de doble cuerda, tonos armónicos, tonos resultantes o
261

armónicos a bassa, brillantez en la velocidad, calidad humana del tono.
[La nobleza y la fluidez del tono resultan muy atractivas para figurar en esta lista, pero, dado que la nobleza del tono depende del tono y del volumen, y la fluidez del tono depende en parte del tono y en parte del brillo, no se les otorga un lugar como cualidades tonales individuales.]
De las doce cualidades tonales individuales mencionadas, solo dos están absolutamente bajo el control del constructor en todo momento: la altura del sonido y el volumen. Por lo tanto, independientemente del valor tonal de la madera, se pueden reproducir indefinidamente violines con la misma altura y volumen. Reproducir las diez cualidades tonales restantes a voluntad y en todo momento es donde reside la gran dificultad. De las restantes, tres están aproximadamente bajo el control del constructor: la intensidad del sonido, la uniformidad de la potencia sonora y la ausencia de armónicos disonantes.
Se ha descubierto que la intensidad del sonido del violín es producto de cuatro factores, por lo tanto:
1. Arqueamiento de las placas.
2. Ubicación y área de las salidas.
3. Estado de las superficies interiores.
4. Acción elástica inherente de la madera de la caja de resonancia.
Es evidente que tres de estos factores están a disposición del constructor en todo momento; pero el cuarto factor no está a disposición en todo momento. La acción elástica inherente de la madera de la tabla de resonancia, siendo un factor...
262

La cantidad costosa opera para vencer la igualdad en la potencia de diferentes violines. Solo en un grado aproximado puede el trabajador experimentado en madera de caja de resonancia predeterminar la calidad tonal de cualquier muestra dada. La calidad tonal de la madera se puede determinar solo por ensayo; pero, el constructor puede asegurar una intensidad considerable al tono del violín mediante el arqueado, por el área y posición de las salidas, y por el estado de las superficies interiores. Incluso con madera de caja de resonancia que carece mucho de poseer una acción de resorte superlativa, es mi observación que se puede asegurar una mayor intensidad de tono con una superficie interior perfectamente lisa que con madera de caja de resonancia que posee una acción de resorte superlativa mientras las superficies interiores permanecen rugosas y alfombradas con fibras de madera, polvo de madera y suciedad. Pero, el aumento del 80 por ciento en la intensidad del tono, como se registró, no debe tomarse como una cantidad de aumento debido únicamente a la superficie interior perfectamente lisa. Este registro se obtuvo con seis violines comunes seleccionados de un stock general, y su potencia de salida promedio se determinó primero sin cambiar sus condiciones; y, el aumento del 80 por ciento en la potencia de transmisión se logró después de que modificadores de tono como la graduación, la barra, el poste, la profundidad de las costillas, los revestimientos, los bloques, el estado de las superficies interiores, el área de salidas, el diapasón, el puente, el barniz y las cuerdas recibieran mi máxima atención. Por lo tanto, el porcentaje de aumento en la intensidad del tono, debido únicamente a la superficie interior perfectamente lisa, no aparece en el registro, ni he realizado experimentos solo con este factor. El valor preciso de
263

Es necesario determinar este factor importante para aumentar la intensidad del sonido del violín, debido a su interés para los propietarios de violines antiguos que tienen cualidades tonales correctas en todos los sentidos, excepto en la calidad de la intensidad.
Aunque eliminar el polvo de las superficies internas del violín con productos como maíz, trigo o avena es mejor que no hacerlo en absoluto, hay propietarios de violines que ni siquiera se atreven a permitir la entrada de estos inofensivos elementos en ese sagrado interior. Por lo que he podido observar, este temor se debe a la expectativa de que la calidad del sonido se vea afectada por el ruido. Es cierto que existen casos en los que este temor está justificado. Así, cuando un violín presenta un sonido ruidoso, colocar una alfombra sobre sus superficies internas contribuye a "mejorar" el sonido; y, dado que todo aquello que disminuye la potencia del sonido puede aniquilar el ruido, este fenómeno se debe a que la onda de ruido es inherentemente más débil que la onda musical; por lo tanto, la onda de ruido es la primera en desaparecer.
La superficie interior, permanente y perfectamente lisa, solo resulta atractiva para aquellos violinistas que buscan una mayor intensidad de sonido. Para ellos, la filosofía que subyace a esta cuestión tiene su interés; y a ellos me dirijo ahora.
Al discutir los principios de la filosofía relacionados con las superficies interiores del violín, es interesante notar ciertos hechos relacionados con el aire mismo, así: En todos los países, los filósofos coinciden en que el aire, sur-
264

que rodean la tierra, existen como gases; que tales gases tienen forma de esferas; que cada esfera es una molécula; que tales moléculas son demasiado pequeñas para ser medidas; que se tocan entre sí como disparos;
que son compresibles; que cuando se liberan de la compresión, retoman la forma esférica con energía; que la fuerza se comunica de molécula a molécula por su expansión y contracción; que, en el aire no confinado, la energía se dispersa por igual en todas las direcciones; que tales líneas de dispersión pueden concentrarse mediante paredes confinantes; que la concentración de tales líneas aumenta la distancia recorrida por la energía en el golpe original; que tal energía puede reflejarse en superficies sólidas y lisas sin pérdida de fuerza; que tal energía no solo se detiene, sino que puede aniquilarse totalmente al golpear cuerpos blandos; que la fuerza del golpe original viaja en ángulo recto con la superficie de impacto; que tal línea de viaje, al golpear una superficie sólida y lisa, se refleja en un ángulo igual al ángulo de incidencia; que el volumen del sonido es proporcional al número de moléculas afectadas por el agente de impacto junto con la cantidad de fuerza en sus golpes; que la intensidad del tono (poder de transmisión del tono) es proporcional a la fuerza del golpe y al grado de concentración dado a las líneas de viaje de la onda sonora; que la distancia recorrida por las ondas sonoras puede verse disminuida por las condiciones de los medios reflectantes y por las condiciones meteóricas.
Se omiten otras cualidades del aire que no vienen al caso. Cada uno de estos aspectos relacionados con el aire está directamente vinculado a la producción del sonido del violín.
265

El tamaño infinitamente pequeño de la molécula de aire posee un gran interés para el estudiante de violín. Este interés se centra en el hecho de que las superficies internas no protegidas del violín no pueden convertirse en medios reflectantes perfectos para las líneas de movimiento de las ondas sonoras. Bajo la lupa, y en el límite del pulido, estas superficies aún presentan crestas, valles y bocas abiertas de cavernas. Es evidente que cada una de esas crestas y valles opera para desviar abruptamente las líneas de movimiento de las ondas sonoras; en otras palabras, para anular la igualdad en los ángulos de incidencia y reflexión; que esas cavernas abiertas operan para aniquilar el movimiento de las ondas sonoras; que esas fuerzas desintegradoras, el calor y la humedad, operan para agudizar esas crestas, para profundizar esos valles y para aumentar la capacidad de esas cavernas; que todos esos agentes que actúan sobre las superficies internas del violín, se combinan para disminuir la concentración del movimiento de las ondas sonoras en las salidas; que tal disminución anula la intensidad del tono; que todo este efecto sobre el tono se debe al tamaño infinitamente pequeño de la molécula de aire y a las superficies infinitamente rugosas y desprotegidas del interior del violín.
Desde la superficie interior, permanente y perfectamente lisa, se aprecia claramente la relación entre la potencia y la intensidad del sonido del violín.
También resulta evidente que, para proteger dicha superficie, es necesario recurrir a un agente protector que permita obtener un alto grado de pulido.
Aparte de la intensidad del tono, existen, en este sentido, otros efectos que resultan de interés para los usuarios de violín que desean la máxima potencia tonal.
266

conjunción con efectos tales como los producidos por:
1. Tiempos de reflexión de las ondas sonoras antes de llegar a las salidas.
2. El sólido, o que cede como medio reflectante.
3. Estado de la superficie de golpeo.
He observado que el número de veces que las ondas sonoras se reflejan antes de llegar a las salidas influye notablemente tanto en el brillo del sonido (es decir, la nitidez del tono en la rapidez de sucesión) como en su potencia. Obviamente, los tiempos de reflexión pueden ser tan grandes que disminuyen la fuerza de cualquier proyectil; asimismo, es evidente que los tiempos de reflexión retrasan la llegada de cualquier proyectil a un punto determinado. En el violín, estos principios se demuestran mediante el método de graduación de la caja de resonancia. Así, cuando la reducción del grosor sitúa el punto de mayor amplitud de oscilación cerca de los extremos de la placa, tanto el brillo como la potencia del sonido disminuyen. Es evidente que la mayor fuerza de los golpes entrantes de la placa se encuentra en el punto de mayor amplitud de oscilación; asimismo, que los tiempos de reflexión de las ondas sonoras aumentan con la distancia desde el punto de origen hasta las salidas. Los experimentos repetidos con frecuencia para localizar el punto de oscilación más amplia me han llevado a la conclusión de que, con este método de graduación de la caja de resonancia, la máxima potencia del tono es imposible. En tales experimentos, se demostró de manera concluyente que colocar el punto de oscilación más amplia a la mitad de la posición de
El puente hacia los extremos de la placa produjo el tono máximo.
267

potencia. Desde mi punto de vista, el aumento de la potencia tonal tras este último método de graduación se debe a dos hechos: la oscilación más amplia posible para las fibras de la caja de resonancia solo se puede asegurar en un punto intermedio entre el puente y los extremos de la placa; y, al acortar la distancia desde el punto de origen hasta las salidas, disminuyendo los tiempos de reflexión, se reduce tanto el retardo como la pérdida de fuerza de las ondas sonoras que reciben el mayor impacto; por lo tanto, mayor brillo del tono y mayor potencia tonal.
Es fácil demostrar en papel que, a medida que aumenta la curvatura, disminuyen los tiempos de reflexión de las ondas sonoras. De hecho, esta proposición es evidente por sí misma: en el violín de caja, las placas son paralelas; por lo tanto, las ondas sonoras que se originan en la placa superior y viajan en ángulo recto con respecto al agente percutor, según la ley, deben tocar la parte posterior en un punto perpendicular al punto de origen; por consiguiente, la onda reflejada debe viajar de regreso a la placa superior directamente sobre la línea de incidencia; por lo tanto, no hay progresión del movimiento de las ondas sonoras hacia las salidas. Si solo se le da a la placa superior una ligera curvatura, entonces las ondas sonoras que se originan en ella no tocarán la parte posterior en un punto perpendicular al punto de origen, sino que incidirán en un punto más cercano a las salidas y se reflejarán en un ángulo igual al ángulo de incidencia según la ley. Así, el movimiento de las ondas sonoras, en cada reflexión, se acerca a las salidas; pero, en este caso, solo llega a las salidas después de muchas reflexiones. Si ahora la parte posterior
268

Si se le da una curvatura de 1-8, las reflexiones se reducen en 1-2. Es evidente, por lo tanto, que las reflexiones disminuyen con la altura de la curvatura; pero no se deduce que la potencia del sonido siga un grado indefinido de curvatura, ya que aumentar la altura de la curvatura aumenta la resistencia; por lo tanto, la curvatura de las placas del violín puede ser tan grande que la fuerza en las cuerdas no puede superar dicha resistencia.
La experiencia, y solo la experiencia, determina que cuando la oscilación más amplia de las fibras de la caja de resonancia se coloca cerca de los extremos de la placa, el arqueamiento, en la posición del puente, no debe ser menor que t para asegurar un brillo de tono satisfactorio; y, cuando el punto de oscilación más amplia se coloca a la mitad desde la posición del puente hasta los extremos de la placa, entonces el arqueamiento no debe ser menor que 1 pulgada. Se demostró previamente que cuanto menor es la altura del arqueamiento, mayor es
la susceptibilidad a la fuerza, por lo tanto altura 7
de
La curvatura se convierte en un factor clave para lograr la máxima potencia tonal.
Dado que la caja de resonancia golpea el aire contenido para producir sonido, su superficie de golpeo tiene interés en este sentido. Es evidente que la fuerza de cualquier golpe puede disminuirse mediante una alfombra suave sobre la superficie del medio de golpeo. Así, cuando la superficie interior de la caja de resonancia está cubierta con una alfombra de finas fibras de madera, la fuerza de sus golpes sobre el aire contenido disminuye en proporción al grosor de dicha alfombra; o, dicho de otro modo, la fuerza de sus golpes aumenta por la alfombra permanente.
269

y una superficie interior perfectamente lisa. En este sentido, la placa trasera posee interés más allá del estado de su superficie interior. Como se indicó anteriormente, puedo asegurar una mayor intensidad de tono al tratar la parte trasera únicamente como un medio reflectante. Por lo tanto, la rigidez en la parte trasera, suficiente para resistir firmemente la fuerza de las ondas sonoras de carga, se convierte en una característica valiosa. Pero, si la parte trasera fuerte se mantiene o no inmóvil en el
Desconozco la presencia de dicha carga. Solo sé que el temblor violento de la parte posterior reduce la intensidad del tono, como se demuestra fácilmente con la prueba de larga distancia al aire libre. He aquí un hecho de doble interés: al realizar la prueba de larga distancia para medir la potencia de proyección, se emplea una fuerte presión del arco; por lo tanto, la parte posterior se ve sometida al límite de fuerza dentro de esa caja de resonancia en particular; y, al ceder la parte posterior a tal cantidad de fuerza, se reduce el retroceso de las moléculas de aire. Es evidente que dicho retroceso disminuye a medida que la parte posterior cede, y como la cesión de la parte posterior disminuye con la menor fuerza en la carga de las ondas sonoras, por lo tanto, con menor presión del arco, dicho violín muestra mayor potencia de proyección que con la mayor presión del arco. Este fenómeno puede observarse con frecuencia. En lo que respecta a la asistencia en un conjunto orquestal, dicho violín es inútil; e inútil porque su tono queda ahogado por las ondas sonoras de los instrumentos de armonía.
Muchas veces he demostrado el hecho de que reemplazar esas espaldas debilitadas por otras que poseen 270

La rigidez contribuye a aumentar la intensidad del sonido. Me parece claro que la rigidez en la parte posterior debe ser igual al límite de fuerza en la caja de resonancia, incluso cuando esta alcanza su máxima oscilación por la acción de las cuerdas al aire bajo la máxima presión del arco; y, porque solo así se puede emplear el límite de presión del arco sin perjudicar la potencia de proyección.
Demostraré la razón de esta conclusión mediante la acción de esta pelota elástica al rebotar contra aquella pared de ladrillos y contra aquella pared divisoria de madera ligera. Al lanzar la pelota contra la pared de ladrillos, la distancia que recorre es proporcional a la fuerza aplicada en el instante del impacto; por lo tanto, cuanto mayor sea la fuerza, mayor será el retroceso. Sin embargo, contra la pared divisoria de madera ligera el resultado es muy diferente. Comenzando con una fuerza moderada aplicada a la pelota y aumentándola gradualmente hasta que iguale la rigidez de la madera, la distancia de retroceso aumenta con la cantidad de fuerza; pero, al aplicar a la pelota una fuerza que excede la rigidez de la madera ligera, la distancia de retroceso disminuye considerablemente. Es evidente que la deformación de la pared le quita a la pelota su energía elástica. Las moléculas de aire son elásticas; y la deformación de la pared les impide mostrar su energía elástica; por lo tanto, disminuye su capacidad de transporte. Pero, si la rigidez de la pared fuera exactamente igual a la mayor fuerza aplicada a la pelota, entonces la distancia de retroceso habría superado las distancias anteriores, porque la energía elástica en la madera, sumada a la energía elástica
en la pelota funciona para aumentar la distancia
271

del retroceso. Por la acción del muro divisorio, es evidente que tal gran aumento en la distancia de retroceso solo es posible con una cierta cantidad de fuerza en la bola; una cantidad menor de fuerza no puede despertar el mismo grado de energía en la madera; mientras que una cantidad mayor de fuerza en la bola provoca la deformación de la madera con sus desastrosos resultados.
Este es un punto de gran interés para el estudiante de violín.
Se trata de la vieja cuestión de regular la potencia sonora del violín mediante el trabajo en la tapa trasera. Las siguientes conclusiones, aunque se basan en repetidas demostraciones prácticas, se presentan únicamente como las conclusiones de una persona; y tales conclusiones pueden tener que sostenerse sin el apoyo de otros. He observado que la rigidez de la tapa trasera puede determinarse de tal manera que aumente los tonos fundamentales o abiertos de cada cuerda; pero, en ningún caso he observado un aumento igual para ningún otro tono. La razón de este fenómeno parece estar claramente indicada por la acción de la bola y la pared divisoria ligera. Así: Cuando la fuerza en la bola es exactamente igual a la acción del resorte en la pared, entonces el retroceso de la bola alcanza el máximo, pero, con menos fuerza en la bola, el retroceso disminuye. Aplicado al violín de la siguiente manera: El desarrollo de la mayor fuerza en las cuerdas exige el uso de toda su longitud; por lo tanto, acortar las cuerdas opera para disminuir la fuerza; por lo tanto, cuando la rigidez de la tapa trasera es exactamente igual a la fuerza en toda la longitud de las cuerdas, se deduce que los tonos producidos al acortar las cuerdas no pueden ser igualmente
272

Aumentado. Es evidente que un violín que posee tonos abiertos potentes, mientras que una marcada debilidad tonal se produce al acortar sucesivamente las cuerdas, tiene un valor tonal insignificante. Por lo tanto, la máxima potencia tonal del violín y la perfecta suavidad de sus superficies internas están inseparablemente ligadas; y la perfecta suavidad de la madera del violín es imposible.
